% chapters/cover.tex — Make LLM-basic 표지
% 레트로 테크니컬 (블루프린트/회로도 스타일) 단색 일러스트
% 한 페이지에 정확히 맞도록 설계

\begin{titlepage}
\thispagestyle{empty}
\setlength{\parindent}{0pt}

% ============ 배경 + 프레임 (절대 위치) ============
\begin{tikzpicture}[remember picture, overlay]
  % 배경 (크림색 종이 느낌)
  \fill[paperbg] (current page.south west) rectangle (current page.north east);

  % 바깥 프레임 (이중선)
  \draw[inkblue, line width=1.8pt]
    ($(current page.north west)+(1.2cm,-1.2cm)$)
    rectangle
    ($(current page.south east)+(-1.2cm,1.2cm)$);
  \draw[inkblue, line width=0.5pt]
    ($(current page.north west)+(1.4cm,-1.4cm)$)
    rectangle
    ($(current page.south east)+(-1.4cm,1.4cm)$);

  % 상단: 회로도 패턴 (신경망 노드 + 연결선)
  \begin{scope}[shift={($(current page.north)+(0,-3.5)$)}, inkblue, line width=0.5pt]
    \foreach \x in {-4.5, -1.5, 1.5, 4.5} {
      \filldraw[inkblue] (\x, 0) circle (3pt);
    }
    \foreach \x in {-5, -2.5, 0, 2.5, 5} {
      \filldraw[inkblue] (\x, -1.2) circle (3pt);
    }
    \foreach \i in {-4.5, -1.5, 1.5, 4.5} {
      \foreach \j in {-5, -2.5, 0, 2.5, 5} {
        \draw[inkblue!25, line width=0.2pt] (\i, 0) -- (\j, -1.2);
      }
    }
    \filldraw[inkblue] (0, -2.4) circle (4pt);
    \draw[inkblue!60] (0, -1.2) -- (0, -2.4);
  \end{scope}

  % 하단: 트랜스포머 블록 미니 다이어그램
  \begin{scope}[shift={($(current page.south)+(0,4.5)$)}, inkblue, line width=0.5pt]
    \draw[inkblue, fill=inkblue!5, rounded corners=2pt] (-3, 0) rectangle (3, 0.7);
    \node[inkblue, font=\small\sffamily\bfseries] at (0, 0.35) {Transformer Block};
    \draw[inkblue, fill=inkblue!10, rounded corners=2pt] (-2.4, 0.9) rectangle (2.4, 1.4);
    \node[inkblue, font=\scriptsize\sffamily] at (0, 1.15) {Multi-Head Attention};
    \draw[inkblue, fill=inkblue!10, rounded corners=2pt] (-2.4, -0.7) rectangle (2.4, -0.2);
    \node[inkblue, font=\scriptsize\sffamily] at (0, -0.45) {Feed-Forward};
    \draw[-{Stealth[length=4pt]}, inkblue!70] (0, 0.7) -- (0, 0.9);
    \draw[-{Stealth[length=4pt]}, inkblue!70] (0, -0.2) -- (0, 0);
  \end{scope}

  % 좌측: 측정 눈금자 (레트로 테크니컬 느낌)
  \begin{scope}[shift={($(current page.west)+(1.8,2)$)}, inkblue!50, line width=0.3pt]
    \foreach \y in {0, 0.5, 1, 1.5, ..., 18} {
      \draw ($(0, -\y)$) -- ($(0.3, -\y)$);
    }
    \draw[inkblue, line width=0.5pt] (0, 0) -- (0, -18);
  \end{scope}
\end{tikzpicture}

% ============ 텍스트 콘텐츠 (상단부터 흐름) ============
\vspace*{0.3\textheight}

\begin{center}
% 메인 제목
{\fontsize{56pt}{68pt}\selectfont\bfseries\sffamily\color{inkblue} Make LLM}\\[0.3cm]

% 부제
{\fontsize{32pt}{38pt}\selectfont\bfseries\sffamily\color{inkgray} basic}\\[1.0cm]

% 구분선
\begin{tikzpicture}
  \draw[inkblue, line width=0.8pt] (0,0) -- (9,0);
\end{tikzpicture}\\[0.4cm]

% 설명 문구
{\large\itshape\sffamily\color{inkgray} 토크나이저부터 트랜스포머까지\\[0.2cm]
처음부터 직접 만드는 대형 언어 모델}\\[0.5cm]

\begin{tikzpicture}
  \draw[inkblue, line width=0.8pt] (0,0) -- (9,0);
\end{tikzpicture}\\[0.6cm]

% 저자 / 출판사 (표지 하단)
{\large\sffamily\bfseries\color{inkblue} dirmich}\\[0.2cm]
{\normalsize\sffamily\color{inkgray} Highmaru Press}
\end{center}

\end{titlepage}

% 표지 다음 페이지로 강제 이동 (섞임 방지)
\clearpage
